\chapter{MISE EN CONTEXTE}
\label{chap:chap_1}
%\section{introduction}
\listoftodos[Notes]%ceci permet d'afficher la liste dee commentaires du professeur

\section*{Introduction}

Les progrès récents ...

\section{Contexte de la recherche}
Le contexte de la recherche doit être présenté de manière à fournir une compréhension approfondie du sujet étudié. Il doit inclure une revue de la littérature pertinente, en mettant en évidence les travaux antérieurs et les lacunes dans la recherche existante. Le contexte doit également expliquer l'importance du sujet de recherche et justifier la nécessité de mener cette étude. Il est important de présenter le contexte de manière claire et concise, en évitant les détails inutiles et en se concentrant sur les aspects les plus pertinents pour la recherche. Le contexte doit être présenté de manière logique, en utilisant des transitions pour relier les différentes idées et en fournissant des exemples concrets pour illustrer les points clés. Il est également important de s'assurer que le contexte est aligné avec les objectifs de recherche et la problématique posée dans l'introduction, en fournissant une base solide pour la suite du document.
La revue de la littérature présenté dans cette section est courte et directement centré sur le projet de recherche de manière à présenter les travaux les plus pertinents pour le projet de recherche et permette l'identification des lacunes dans la recherche existante et être en mesure d'écrire adéquatement une problématique claire.

\section{Problématique}

La rédaction d'une problématique doit être claire, concise et bien structurée. La problématique annoce le vérou scientifique du projet de recherche, en identifiant les lacunes dans la littérature existante de la section précédente et en présentant les questions de recherche qui seront abordées dans le document. Il est important de formuler la problématique de manière à ce qu'elle soit compréhensible pour les lecteurs, en évitant le jargon technique ou les termes ambigus. La problématique doit être présentée de manière logique, en utilisant l'introduction générale du sujet et de l'analyse des travaux antérieurs et des lacunes précédemment identifiées. La problématique adresse les questions de recherche spécifiques qui seront abordées dans le document. Il est également important de s'assurer que la problématique est alignée avec les objectifs de recherche et qu'elle contribue de manière significative à la compréhension du sujet étudié.

L'introduction  des...
La problématique centrale de ce projet de recherche réside...


\section{Objectifs}
Les objectifs doivent être clairement définis et alignés avec la problématique de recherche. Ils doivent être spécifiques, mesurables, atteignables, pertinents et temporellement définis (SMART). Les objectifs doivent être formulés de manière à guider la recherche et à fournir une base pour l'évaluation des résultats. Ils doivent également être présentés de manière logique, en commençant par les objectifs généraux et en les décomposant en objectifs spécifiques. Il est important de s'assurer que les objectifs sont réalisables dans le cadre du projet de recherche et qu'ils contribuent de manière significative à la compréhension du sujet étudié. 

L'objectif principal de ce projet de recherche est de déterminer...

\subsection{Conseils de rédaction}

Pour enlever le contour de l'image, on peut utiliser les paramètres trim et clip. La commande est :\par

\begin{verbatim}
includegraphics[trim={A cm B cm D cm E cm},clip, width=\linewidth],
\end{verbatim}
avec: \par
trim=A:left B:bottom C:right D:top\par 

Voici un exemple de figure côte à côte:

\begin{figure}[!htbp]
\centering
    \begin{minipage}{0.45\linewidth}
        \centering
        \includegraphics[width=\textwidth]{usine.JPG}
        \caption{Example of the hybrid mobile robot with UR5 serial manipulator for the use case study. Source of 3D files, CC-BY \cite{vanel}}\label{fig:robotUR5}
    \end{minipage}
    \begin{minipage}{0.45\linewidth}
        \centering
        \includegraphics[width=\textwidth]{usine.JPG}
        \caption{3D simulation of an electrical substation on Gazebo in ROS. Source of 3D files, CC-BY \cite{vanel}.}\label{fig:3DSousStation}
    \end{minipage}
\end{figure}

\subsubsection{Utilisation des accronymes}
Pour les définitions des acronymes, on doit les ajouter manuellement, par exemple pour \ac{BMS}.
Le même principe s'applique pour \ac{TSP}.
À chaque fois que l'acronyme est utilisé, on doit prendre la commande \ac{TSP}, après l'avoir défini à la première référence.

\subsubsection{Notes de bas de page}

Vous pouvez utiliser les notes de bas de page pour inclure un lien\footnote{\url{https://www.uqac.ca}} en rapport avec votre texte, ou pour donner plus de précisions\footnote{Cette note de bas de page propose plus d'informations}.

\subsubsection{Exemple de listes}
\label{subsec:subsec_list}

\begin{enumerate}
  \item {\bf premier point (gras) ;}
  \item {\em second point (italique) ;}
  \item {\Large troisième point (gros) ;}
      \begin{enumerate}
          \item {\small premier sous-point en petit}
          \item {\tiny second sous-point (petit)}
          \item {\Huge troisième sous-point (très gros)}
      \end{enumerate}
  \item[$\bullet$] {\sf point avec une puce (sans serif)}
  \item[$\circ$] {\sc point avec un autre style de puce (petites lettres capitales)}
\end{enumerate}


\section*{Conclusion}
\label{sec:sec_examples_chap_1}

La conclusion n'est pas une discussion. Elle doit être structurée de manière à résumer les points clés abordés dans le chapitre, à présenter les implications de ces points pour la suite du document et à faire le lien avec les chapitres suivants.

La conclusion doit comprendre une synthèse des points clés abordés dans le chapitre, ainsi que les implications de ces points pour la suite du document. Elle doit également faire le lien avec les chapitres suivants en présentant les questions de recherche qui seront abordées dans les sections suivantes.

La conclusion doit être concise et claire, en évitant de répéter les informations déjà présentées dans le chapitre. Elle doit plutôt se concentrer sur les implications et les perspectives futures de la recherche présentée dans le chapitre.

Il est requis de réaliser un rappel des hypothèses afin de la confirmer ou des les infirmer, ainsi que de présenter les travaux futurs qui seront réalisés dans les sections suivantes.

Ce projet de recherche vise à résoudre...

\newpage
